\section{Introduction}
This article is concerned with dif\/ferential calculus on real or complex
projective space, in\-variant under projective transformations. These
transformations constitute a semisimple Lie group and projective space is hence
a homogeneous space of the form~$G/P$ for~$G$ semisimple. The subgroup~$P$ is
parabolic and, more generally, dif\/ferential geometries modelled on homogeneous
spaces of this form are known as `parabolic'~\cite{thebook}. Projective space
gives rise to projective dif\/ferential geometry in this sense. Conformal and CR
geometry are included amongst other examples of parabolic dif\/ferential
geometry. The interplay between the symmetries of projective space and its
invariant dif\/ferential operators is mediated by representation theory. In this
article we present the `BGG complex' on projective space as perhaps the
simplest of these constructions in one of the simplest of settings. We
anticipate that our approach will extend to $G/P$ in general and perhaps to
other homogeneous spaces. In~\cite{o}, Olver introduced complexes of
dif\/ferential operators on Euclidean space acting between `hyperforms'. These
are BGG complexes constructed directly, employing only af\/f\/ine invariance and
the associated Schur functors. He constructs some `easy' examples and observes
that ``Other examples, of greater complexity, can of course be constructed at
will, but the expressions rapidly get out of hand, even in low dimensional
spaces''. We maintain that the extra symmetry that the operators exhibit under
projective transformations allows one to control these expressions more
ef\/fectively. The fascinating combination of symmetries and dif\/ferential
operators is a def\/ining feature in the work of Willard Miller to whom we
dedicate this article.

The Bernstein--Gelfand--Gelfand (BGG) complexes on ${\mathbb{RP}}_n$ or
${\mathbb{CP}}_n$ are, by now, well-known complexes of vector bundles and
dif\/ferential operators between them generalising the de~Rham complex on
${\mathbb{RP}}_n$ and the holomorphic de~Rham complex on~${\mathbb{CP}}_n$,
respectively. An introduction to such complexes is given in~\cite{eVariations}
and the particular case of projective space is discussed in~\cite{eProjective}.
Usually, their construction involves choosing so-called `splitting
operators'~\cite{cd,css4} constructed from Kostant's Laplacian \cite{k} or
`quabla operator' \cite{cd} or from the Jantzen--Zuckerman translation
functor~\cite{v}. Here, we avoid the direct use of splitting operators, instead
relying only on diagram chasing, as is already done
in~\cite{eVariations,eProjective} in simple cases. In fact, the BGG complex of
holomorphic dif\/ferential operators on ${\mathbb{CP}}_2$ was already constructed
in this manner~\cite[p.~351]{eDuality} before it was realised by John
Rice~\cite{r} that complexes like this were dual to Lepowsky's construction
\cite{l} on the level of Verma modules. (On ${\mathbb{CP}}_1$ there is just one
family of BGG operators, already singled out for their invariance
in~\cite[Proposition~2.1]{et}.) In this article, for simplicity and cleanliness
we employ a spectral sequence to ef\/fect the diagram chasing. We employ
projective invariance to derive explicit formul{\ae} for the operators in the
projective BGG complex on ${\mathbb{RP}}_n$ in terms of the usual round metric
on the sphere.

As is often done in dif\/ferential geometry, when it is necessary to write out
tensors and their natural operations and we shall adorn them with upper or
lower indices corresponding to the tangent or cotangent bundle respectively.
For example, a vector f\/ield can be written as~$X^a$, a~one-form as $\omega_a$,
and the natural pairing between them as $X^a\omega_a$ in accordance with the
`Einstein summation convention'. For any tensor $\phi_{abc}$ we shall write its
symmetric part as $\phi_{(abc)}$ and its skew part as $\phi_{[abc]}$. For
example, to say that $\omega_{ab}$ is a two-form is to say that
\[\omega_{ab}=-\omega_{ba}\qquad
\mbox{or, equivalently,}\qquad\omega_{ab}=\omega_{[ab]}\qquad
\mbox{or, equivalently,}\qquad\omega_{(ab)}=0\]
and then
\[\nabla_{[a}\omega_{bc]}\qquad\mbox{and}\qquad
X^a\nabla_a\omega_{bc}-2(\nabla_{[b}X^a)\omega_{c]a},\]
for any torsion-free connection~$\nabla_a$, are the exterior derivative of
$\omega_{ab}$ and the Lie derivative of $\omega_{ab}$ in the direction of the
vector f\/ield~$X^a$, respectively. Such formul{\ae} are not meant to imply any
choice of local co\"ordinates. More precisely, this is Penrose's `abstract
index notation'~\cite{OT} and it formalises the conventions used by many
classical authors~-- see, for example, the discussion of projective dif\/ferential
geometry by Schouten~\cite{s}.

We shall view ${\mathbb{RP}}_n$ as a homogeneous space
\[{\mathbb{RP}}_n={\mathrm{SL}}(n+1,{\mathbb{R}})/P=G/P,\qquad\mbox{where}\quad
P=\left\{
\left\lgroup\begin{tabular}{c|ccc}$\ast$&$\ast$&$\!\cdots\!$&$\ast$\\ \hline
$0$\\
\raisebox{2pt}[15pt]{$\vdots$}&&
\raisebox{-2pt}[0pt][0pt]{\makebox[0pt]{\Huge$\ast$}}\\
$0$\end{tabular}\right\rgroup\right\}\]
or as a quotient of the `projective sphere'
\[S^n={\mathrm{SL}}(n+1,{\mathbb{R}})/P=G/P,\qquad\mbox{where}\quad
P=\left\{
\left\lgroup\begin{tabular}{c|ccc}$\lambda$&$\ast$&$\!\cdots\!$&$\ast$\\ \hline
$0$\\
\raisebox{2pt}[15pt]{$\vdots$}&&
\raisebox{-2pt}[0pt][0pt]{\makebox[0pt]{\Huge$\ast$}}\\
$0$\end{tabular}\right\rgroup\mbox{ s.t.\ }\lambda>0\right\}\]
under the antipodal map. In either case, it is convenient to write the
associated Lie algebra as
\begin{gather}\label{graded}
{\mathfrak{sl}}(n+1,{\mathbb{R}})={\mathfrak{g}}=
{\mathfrak{g}}_-\oplus{\mathfrak{g}}_0\oplus{\mathfrak{g}}_+,
\end{gather}
where{\samepage
\begin{gather*}
{\mathfrak{g}}_-=\left\{\!
\left\lgroup\!\begin{tabular}{c|ccc}$0$&$0$&$\!\cdots\!$&$0$\\ \hline
$\ast$\\
\raisebox{2pt}[15pt]{$\vdots$}&&
\raisebox{-2pt}[0pt][0pt]{\makebox[0pt]{\Huge$0$}}\\
$\ast$\end{tabular}\!\right\rgroup\!\right\},\!\quad
{\mathfrak{g}}_0=\left\{\!
\left\lgroup\!\begin{tabular}{c|ccc}$\ast$&$0$&$\!\cdots\!$&$0$\\ \hline
$0$\\
\raisebox{2pt}[15pt]{$\vdots$}&&
\raisebox{-2pt}[0pt][0pt]{\makebox[0pt]{\Huge$\ast$}}\\
$0$\end{tabular}\!\right\rgroup\!\right\},\!\quad
{\mathfrak{g}}_+=\left\{\!
\left\lgroup\!\begin{tabular}{c|ccc}$0$&$\ast$&$\!\cdots\!$&$\ast$\\ \hline
$0$\\
\raisebox{2pt}[15pt]{$\vdots$}&&
\raisebox{-2pt}[0pt][0pt]{\makebox[0pt]{\Huge$0$}}\\
$0$\end{tabular}\!\right\rgroup\!\right\}
\end{gather*}
and then ${\mathfrak{p}}={\mathfrak{g}}_0\oplus{\mathfrak{g}}_+$.}

From now on we shall discuss only real projective space ${\mathbb{RP}}_n$ or
its double cover, the sphere~$S^n$. The complex case is completely parallel
with real numbers being replaced by complex numbers everywhere and by working
in the holomorphic category rather than the smooth.

\section{An outline of the construction}\label{outline}

If ${\mathbb{V}}$ is a f\/inite-dimensional representation of~$P$, we shall
denote by $V$ the induced homogeneous vector bundle on~$G/P$ constructed as
\[V=G\times_P{\mathbb{V}}=G\times{\mathbb{V}}/\sim,
\qquad\mbox{where }(g,v)\sim\big(g p,p^{-1}v\big),\quad \forall\, p\in P.\]
Notice that if ${\mathbb{V}}$ is actually a $G$-module restricted to~$P$, then
$V$ is canonically trivialised
\begin{equation}\label{twist}
V=G\times_P{\mathbb{V}}\cong G/P\times{\mathbb{V}}\qquad\mbox{by}\quad
(g,v)\mapsto(gP,g v)\end{equation}
as a vector bundle (but not as a homogeneous vector bundle). Hence, in this
case $V$ is naturally equipped with a $G$-equivariant f\/lat connection~$\nabla$
obtained by transporting to $V$ the exterior derivative
$d:\Lambda^0\otimes{\mathbb{V}}\to\Lambda^1\otimes{\mathbb{V}}$ with
values in~${\mathbb{V}}$. More generally, the coupled de~Rham sequence
\begin{equation}\label{coupled}V\xrightarrow{\,\nabla\,}\Lambda^1\otimes V
\xrightarrow{\,\nabla\,}\Lambda^2\otimes V
\xrightarrow{\,\nabla\,}\Lambda^3\otimes V \xrightarrow{\,\nabla\,}\cdots
\xrightarrow{\,\nabla\,}\Lambda^{n-1}\otimes V
\xrightarrow{\,\nabla\,}\Lambda^n\otimes V\to 0\end{equation} is exact on the
level of germs and provides a resolution of ${\mathbb{V}}$ as a locally
constant sheaf on~$G/P$.

This general reasoning holds on any homogeneous space $G/P$ but the following
discussion is specif\/ic to ${\mathbb{RP}}_n$ or~$S^n$. Suppose ${\mathbb{V}}$ is
irreducible as a $G$-module. In this case we shall see that as a $P$-module
${\mathbb{V}}$ is f\/iltered
\begin{equation}\label{filtration}
{\mathbb{V}}={\mathbb{V}}_0+{\mathbb{V}}_1+{\mathbb{V}}_2+
\cdots+{\mathbb{V}}_{N-1}+{\mathbb{V}}_N\end{equation} meaning that
these are the subquotients listed in a natural order, starting on the
left with the smallest quotient of $\mathbb{V}$.  (In other words
${\mathbb{V}}_N$ is the smallest $P$-submodule in the f\/iltration, the
quotient ${\mathbb{V}}/{\mathbb{V}}_N$ has a f\/iltration
${\mathbb{V}}_0+{\mathbb{V}}_1+{\mathbb{V}}_2+
\cdots+{\mathbb{V}}_{N-1}$, and the meaning is now clear by induction.)
It follows that the bundle $V$ is correspondingly f\/iltered
\begin{equation}\label{correspondingfiltration}
V=V_0+V_1+V_2+\cdots+V_{N-1}+V_N\end{equation}
and now we claim that the connection $\nabla:V\to\Lambda^1\otimes V$, and
consequently the whole complex~(\ref{coupled}), is compatible with this
f\/iltration (as detailed in Theorem~\ref{compatible} below). The spectral
sequence of a~f\/iltered complex~\cite{c} now comes into play, having as its
$E_0$-level the following.
\begin{equation}\label{E0}\raisebox{-70pt}{\begin{picture}(300,148)(0,10)
\put(0,20){\vector(1,0){100}}
\put(0,20){\vector(0,1){80}}
\put(93,14){\makebox(0,0){$p$}}
\put(-4,91){\makebox(0,0){$q$}}
\put(10,150){\makebox(0,0){$V_0$}}
\put(50,150){\makebox(0,0){$\Lambda^1\otimes V_0$}}
\put(100,150){\makebox(0,0){$\Lambda^2\otimes V_0$}}
\put(150,150){\makebox(0,0){$\Lambda^3\otimes V_0$}}
\put(200,150){\makebox(0,0){$\Lambda^4\otimes V_0$}}
\put(250,150){\makebox(0,0){$\Lambda^5\otimes V_0$}}
\put(290,150){\makebox(0,0){$\cdots$}}
\put(290,120){\makebox(0,0){$\cdots$}}
\put(290,90){\makebox(0,0){$\cdots$}}
\put(290,60){\makebox(0,0){$\cdots$}}
\put(290,30){\makebox(0,0){$\cdots$}}
\put(50,135){\makebox(0,0){$\uparrow\scriptstyle\partial$}}
\put(100,135){\makebox(0,0){$\uparrow\scriptstyle\partial$}}
\put(150,135){\makebox(0,0){$\uparrow\scriptstyle\partial$}}
\put(200,135){\makebox(0,0){$\uparrow\scriptstyle\partial$}}
\put(250,135){\makebox(0,0){$\uparrow\scriptstyle\partial$}}
\put(50,120){\makebox(0,0){$V_1$}}
\put(100,120){\makebox(0,0){$\Lambda^1\otimes V_1$}}
\put(150,120){\makebox(0,0){$\Lambda^2\otimes V_1$}}
\put(200,120){\makebox(0,0){$\Lambda^3\otimes V_1$}}
\put(250,120){\makebox(0,0){$\Lambda^4\otimes V_1$}}
\put(100,105){\makebox(0,0){$\uparrow\scriptstyle\partial$}}
\put(150,105){\makebox(0,0){$\uparrow\scriptstyle\partial$}}
\put(200,105){\makebox(0,0){$\uparrow\scriptstyle\partial$}}
\put(250,105){\makebox(0,0){$\uparrow\scriptstyle\partial$}}
\put(100,90){\makebox(0,0){$V_2$}}
\put(150,90){\makebox(0,0){$\Lambda^1\otimes V_2$}}
\put(200,90){\makebox(0,0){$\Lambda^2\otimes V_2$}}
\put(250,90){\makebox(0,0){$\Lambda^3\otimes V_2$}}
\put(150,75){\makebox(0,0){$\uparrow\scriptstyle\partial$}}
\put(200,75){\makebox(0,0){$\uparrow\scriptstyle\partial$}}
\put(250,75){\makebox(0,0){$\uparrow\scriptstyle\partial$}}
\put(150,60){\makebox(0,0){$V_3$}}
\put(200,60){\makebox(0,0){$\Lambda^1\otimes V_3$}}
\put(250,60){\makebox(0,0){$\Lambda^2\otimes V_3$}}
\put(200,45){\makebox(0,0){$\uparrow\scriptstyle\partial$}}
\put(250,45){\makebox(0,0){$\uparrow\scriptstyle\partial$}}
\put(200,30){\makebox(0,0){$V_4$}}
\put(250,30){\makebox(0,0){$\Lambda^1\otimes V_4$}}
\end{picture}}\end{equation}
Here, the precise positioning of the co\"ordinate axes is a matter of
convention. The important property of the $E_0$-level is that the
dif\/ferentials $\partial$ are simply homomorphisms of vector bundles and we
shall show that they are induced by a complex of $G_0$-modules
\[{\mathbb{V}}\stackrel{\partial}{\longrightarrow}
{\mathfrak{g}}_-^*\otimes{\mathbb{V}}\stackrel{\partial}{\longrightarrow}
\Lambda^2{\mathfrak{g}}_-^*\otimes{\mathbb{V}}
\stackrel{\partial}{\longrightarrow}
\Lambda^3{\mathfrak{g}}_-^*\otimes{\mathbb{V}}
\stackrel{\partial}{\longrightarrow}\cdots\stackrel{\partial}{\longrightarrow}
\Lambda^{n-1}{\mathfrak{g}}_-^*\otimes{\mathbb{V}}
\stackrel{\partial}{\longrightarrow}
\Lambda^n{\mathfrak{g}}_-^*\otimes{\mathbb{V}},\]
\newpage

\noindent
where
\[G_0=\left\{
\left\lgroup\begin{tabular}{c|ccc}$\ast$&$0$&$\!\cdots\!$&$0$\\ \hline
$0$\\
\raisebox{2pt}[15pt]{$\vdots$}&&
\raisebox{-2pt}[0pt][0pt]{\makebox[0pt]{\Huge$\ast$}}\\
$0$\end{tabular}\right\rgroup\in{\mathrm{SL}}(n+1,{\mathbb{R}})\right\}.\]
Furthermore, we shall show that this complex def\/ines the Lie algebra
cohomology $H^r({\mathfrak{g}}_-,{\mathbb{V}})$, which in turn has been
computed by Kostant~\cite{k}. It follows that the $E_1$-level of the spectral
sequence is rather sparse, typically
\begin{equation}\label{sparse}\raisebox{-70pt}{\begin{picture}(300,148)(0,10)
\put(0,20){\vector(1,0){100}}
\put(0,20){\vector(0,1){80}}
\put(93,14){\makebox(0,0){$p$}}
\put(-4,91){\makebox(0,0){$q$}}
\put(10,150){\makebox(0,0){$H^0$}}
\put(50,150){\makebox(0,0){$0$}}
\put(100,150){\makebox(0,0){$0$}}
\put(150,150){\makebox(0,0){$0$}}
\put(200,150){\makebox(0,0){$0$}}
\put(250,150){\makebox(0,0){$0$}}
\put(290,150){\makebox(0,0){$\cdots$}}
\put(290,120){\makebox(0,0){$\cdots$}}
\put(290,90){\makebox(0,0){$\cdots$}}
\put(290,60){\makebox(0,0){$\cdots$}}
\put(290,30){\makebox(0,0){$\cdots$}}
\put(50,120){\makebox(0,0){$0$}}
\put(100,120){\makebox(0,0){$H^1$}}
\put(150,120){\makebox(0,0){$H^2$}}
\put(200,120){\makebox(0,0){$0$}}
\put(250,120){\makebox(0,0){$0$}}
\put(100,90){\makebox(0,0){$0$}}
\put(150,90){\makebox(0,0){$0$}}
\put(200,90){\makebox(0,0){$0$}}
\put(250,90){\makebox(0,0){$H^3$}}
\put(150,60){\makebox(0,0){$0$}}
\put(200,60){\makebox(0,0){$0$}}
\put(250,60){\makebox(0,0){$0$}}
\put(200,30){\makebox(0,0){$0$}}
\put(250,30){\makebox(0,0){$0$}}
\end{picture}}\end{equation}
where $H^0=V_0$ and, in particular, there is precisely one irreducible bundle
in each diagonal~$E_1^{p,d-p}$ for~$d$ f\/ixed. Since~(\ref{coupled})
resolves~${\mathbb{V}}$, we know that this spectral sequence is also converging
to ${\mathbb{V}}$ and the only way that this can happen is if the dif\/ferentials
f\/it together as a~resolution
\[0\to{\mathbb{V}}
\to H^0\to H^1\to H^2\to H^3\to\cdots\to H^{n-1}\to H^n\to 0.\]
This is the required BGG resolution. The rest of the article is devoted to
f\/illing in the details of this argument.

\section[The filtering of ${\mathbb{V}}$ as a $P$-module]{The f\/iltering of $\boldsymbol{{\mathbb{V}}}$ as a $\boldsymbol{P}$-module}\label{filtering}

Recall the decomposition (\ref{graded}) of
${\mathfrak{g}}={\mathfrak{sl}}(n+1,{\mathbb{R}})$ and consider the element
\[H=\frac{1}{n+1}
\left\lgroup\begin{tabular}{c|ccc}$n$&$0$&$\!\cdots\!$&$0$\\ \hline
$0$&$-1$&&$0$\\
\raisebox{2pt}[15pt]{$\vdots$}&&
\raisebox{1pt}[0pt][0pt]{\makebox[0pt]{\Large$\ddots$}}\\
$0$&$0$&&$-1$\end{tabular}\right\rgroup\in{\mathfrak{g}}_0.\]
uniquely characterised as lying in the centre of ${\mathfrak{g}}_0$ with
$[H,X]=X$, for $X\in{\mathfrak{g}}_+$. It is called the {\em grading
element\/}~\cite{thebook} of the $|1|$-graded Lie algebra~(\ref{graded}). If
the $G$-module ${\mathbb{V}}$ is restricted to~$G_0$, then it splits into
eigenspaces under~$H$. For the standard representation by matrix multiplication
on column vectors for example,
\begin{equation}\label{Hstandard}
H\left\lgroup\begin{array}{c}0\\ v_1\\ \vdots\\ v_n\end{array}\right\rgroup=
-\frac{1}{n+1}
\left\lgroup\begin{array}{c}0\\ v_1\\ \vdots\\ v_n\end{array}\right\rgroup
\qquad\mbox{and}\qquad
H\left\lgroup\begin{array}{c}x\\ 0\\ \vdots\\ 0\end{array}\right\rgroup=
\frac{n}{n+1}
\left\lgroup\begin{array}{c}x\\ 0\\ \vdots\\ 0\end{array}\right\rgroup.
\end{equation}
Rather than use the actual eigenvalues, which are rational in general, let us
subtract the lowest eigenvalue and write
\begin{equation}\label{splits}
{\mathbb{V}}={\mathbb{V}}_0\oplus{\mathbb{V}}_1\oplus{\mathbb{V}}_2
\oplus\cdots\oplus{\mathbb{V}}_{N-1}\oplus{\mathbb{V}}_N\end{equation}
for the eigenspace decomposition, noting that ${\mathfrak{g}}_+$ acts by
${\mathbb{V}}_j\to{\mathbb{V}}_{j+1}$ for all~$j$. It follows that
\begin{equation}\label{submodules}
{\mathbb{V}}^j\equiv{\mathbb{V}}_j\oplus{\mathbb{V}}_{j+1}\oplus\cdots\oplus
{\mathbb{V}}_N\end{equation}
are $P$-submodules of ${\mathbb{V}}$ for all $j$ and we have our
f\/iltration~(\ref{filtration}).

\section[The filtered complex $\Lambda^\bullet\otimes V$ and its spectral sequence]{The f\/iltered complex $\boldsymbol{\Lambda^\bullet\otimes V}$ and its spectral sequence}

The f\/iltration (\ref{submodules}) of ${\mathbb{V}}$ as a $P$-module certainly
induces a f\/iltration
\[
V=V^0\supseteq V^1\supseteq V^2\supseteq\cdots\supseteq V^{N-1}\supset V^N
\]
of $V$ by $G$-homogeneous vector bundles on~$G/P$.
\begin{theorem}\label{compatible} The connection
$\nabla:V\to\Lambda^1\otimes V$ is compatible with this filtration in the
sense that there is a commutative diagram
\[\begin{array}{ccc}V&\stackrel{\nabla}{\longrightarrow}&\Lambda^1\otimes V\\
\rule[.7pt]{.4pt}{6pt}\hspace*{.7pt}\cup
&&\rule[.7pt]{.4pt}{6pt}\hspace*{.7pt}\cup\\
V^k&\stackrel{\nabla}{\longrightarrow}&\Lambda^1\otimes V^{k-1}
\end{array}\]
for all $k=1,2,\ldots,N$.
\end{theorem}
\begin{proof}
Because the assertion is local, without loss of generality it suf\/f\/ices to prove
it on the standard af\/f\/ine co\"ordinate patch, namely
\begin{gather}\label{Q}
{\mathbb{R}}^n=Q/G_0\subset G/P={\mathbb{RP}}_n,\qquad\mbox{where}\quad
Q=\left\{
\left\lgroup\begin{tabular}{c|ccc}$\ast$&$0$&$\!\cdots\!$&$0$\\ \hline
$\ast$\\
\raisebox{2pt}[15pt]{$\vdots$}&&
\raisebox{-2pt}[0pt][0pt]{\makebox[0pt]{\Huge$\ast$}}\\
$\ast$\end{tabular}\right\rgroup\in{\mathrm{SL}}(n+1,{\mathbb{R}})\right\}.\!\!\!\!
\end{gather}
Recall (\ref{splits}) that ${\mathbb{V}}$ splits as a $G_0$-module. Hence the
same is true of $V$ restricted to this patch:%--
\begin{equation}\label{Vsplitting}V|_{{\mathbb{R}}^n}
=V_0\oplus V_1\oplus V_2\oplus\cdots\oplus V_{N-1}\oplus V_N.\end{equation}
To proceed we need a formula for~$\nabla$. The appendix discusses various
natural constructions on a general Lie group $G$, which we now specialise to
be the Abelian Lie group
\[{\mathbb{R}}^n=
G_-=\left\{
\left\lgroup\begin{tabular}{c|ccc}$1$&$0$&$\!\cdots\!$&$0$\\ \hline
$\ast$\\
\raisebox{2pt}[15pt]{$\vdots$}&&
\raisebox{-2pt}[0pt][0pt]{\makebox[0pt]{\LARGE${\mathrm{Id}}$}}\\
$\ast$\end{tabular}\right\rgroup\in{\mathrm{SL}}(n+1,{\mathbb{R}})\right\}.\]
and according to (\ref{nabla}) we f\/ind that $\nabla=d+\theta$, where
\begin{itemize}\itemsep=0pt
\item $V|_{{\mathbb{R}}^n}$ is trivialised by
$\begin{array}{ccc}G_-\times{\mathbb{V}}&
\stackrel{\simeq\:\:}{\rightarrow}&Q\times_{G_0}{\mathbb{V}}\\
\cap&\raisebox{-2pt}{$\nearrow$}&\|\\
Q\times{\mathbb{V}}&&V|_{{\mathbb{R}}^n}
\end{array}$ to def\/ine
$d:V|_{{\mathbb{R}}^n}\to\Lambda^1\otimes V|_{{\mathbb{R}}^n}$,
\item $\theta:V|_{{\mathbb{R}}^n}\to\Lambda^1\otimes V|_{{\mathbb{R}}^n}$ is
def\/ined by the same trivialisation together with the Maurer--Cartan form
$\theta$ on $G_-$; more specif\/ically,
\[V|_{{\mathbb{R}}^n}\cong
G_-\times{\mathbb{V}}\xrightarrow{\,\theta\otimes{\mathrm{Id}}\,}
\Lambda^1\otimes{\mathfrak{g}}_-\otimes{\mathbb{V}}
\xrightarrow{\,{\mathrm{Id}}\otimes\rho\,}\Lambda^1\otimes{\mathbb{V}}\cong
\Lambda^1\otimes V|_{{\mathbb{R}}^n},\]
where $\rho:{\mathfrak{g}}_-\otimes{\mathbb{V}}\to{\mathbb{V}}$ is the
representation of ${\mathfrak{g}}$ on ${\mathbb{V}}$ restricted
to~${\mathfrak{g}}_-$.
\end{itemize}
Evidently, $d$ preserves the splitting (\ref{Vsplitting}) whilst $\theta$
sends $V_k$ to $V_{k-1}$ for all $k=1,2,\ldots, N$. In particular, $\nabla$
sends $V^k=V_k\oplus\cdots$ to
$\Lambda^1\otimes V^{k-1}=\Lambda^1\otimes V_{k-1}\oplus\cdots$, as required.
\end{proof}

\begin{corollary}\label{easypeasy}
The complex~\eqref{coupled} is compatible with the
filtration~\eqref{correspondingfiltration}, i.e.~$\nabla$ sends
$\Lambda^p\otimes V^k$ to $\Lambda^{p+1}\otimes V^{k-1}$.
\end{corollary}

\begin{proof} In fact, in the trivialisation $V|_{{\mathbb{R}}^n}\cong
G_-\times{\mathbb{V}}$ employed in the proof of Theorem~\ref{compatible}
\begin{gather*}
\Lambda^p\otimes V|_{{\mathbb{R}}^n}\cong\Lambda^p\otimes{\mathbb{V}}\ni
\omega\otimes v\stackrel{\nabla}{\longmapsto}
d\omega\otimes v+(-1)^p\omega\wedge(\theta\intprod\rho)v\in
\Lambda^{p+1}\otimes{\mathbb{V}}\cong\Lambda^{p+1}\otimes V|_{{\mathbb{R}}^n}
\end{gather*}
and the conclusion is manifest.\end{proof}

According to this corollary, we may now consider the spectral sequence of the
f\/iltered complex $\nabla:\Lambda^\bullet\otimes V$ on~$G/P$, the $E_0$-level
of which is
\begin{equation}\label{E0pq}
E_0^{p,q}=\Lambda^{p+q}\otimes V_{-q}\qquad\mbox{with dif\/ferential}\quad
\partial: \ E_0^{p,q}\to E_0^{p,q+1},\end{equation}
where we have chosen to normalise (\ref{E0}) by placing $V_0$ at the origin. By
construction, the bundle~$V_k$ on $G/P$ is the homogeneous bundle induced from
${\mathbb{V}}_k={\mathbb{V}}^k/{\mathbb{V}}^{k+1}$ as a $P$-module
(cf.~(\ref{submodules})). We already know that, as an eigenspace for the
grading element~$H$ in the centre of~${\mathfrak{g}}_0$, the vector space
${\mathbb{V}}_k$ is a~$G_0$-module. By regarding it as the quotient
${\mathbb{V}}^k/{\mathbb{V}}^{k+1}$ we are equivalently making~${\mathbb{V}}_k$
into a $P$-module by decreeing that $G_+$ act trivially. By construction, the
$E_0$-dif\/ferential is $G$-equivariant. Furthermore, its def\/inition
\[\begin{array}{@{}ccccccccc}
0&\to&\Lambda^{p+1}\otimes V^{k-2}&\to&\Lambda^{p+1}\otimes V^{k-1}&\to&
\Lambda^{p+1}\otimes V_{k-1}&\to&0\\
&&\nabla\uparrow\phantom{\nabla}&&\nabla\uparrow\phantom{\nabla}&&
\partial\uparrow\phantom{\partial}\\
0&\to&\Lambda^p\otimes V^{k-1}&\to&\Lambda^p\otimes V^k&\to&
\Lambda^p\otimes V_k&\to&0
\end{array}\]
and the Leibniz rule ensure that it is linear over the functions. In other
words $\partial$ is a $G$-equivariant homomorphism of homogeneous bundles and,
as such, must be induced by a homomorphism of $P$-modules
$\Lambda^p({\mathfrak{g}}/{\mathfrak{p}})^*\otimes{\mathbb{V}}_k\to
\Lambda^{p+1}({\mathfrak{g}}/{\mathfrak{p}})^*\otimes{\mathbb{V}}_{k-1}$, which
we shall also denote by~$\partial$. In fact, from the formula for $\nabla$
displayed in the proof of Corollary~\ref{easypeasy}, we see that $\partial$ is
induced by
\begin{equation}\label{E0differential}
\Lambda^p({\mathfrak{g}}/{\mathfrak{p}})^*\otimes{\mathbb{V}}=
\Lambda^p{\mathfrak{g}}_-^*\otimes{\mathbb{V}}
\ni v_{\beta\gamma\cdots\delta}\stackrel{\partial}{\longmapsto}
\rho_{[\alpha}v_{\beta\gamma\cdots\delta]}\in
\Lambda^{p+1}{\mathfrak{g}}_-^*\otimes{\mathbb{V}}=
\Lambda^{p+1}({\mathfrak{g}}/{\mathfrak{p}})^*\otimes{\mathbb{V}},
\end{equation}
where ${\mathbb{V}}$ is regarded as a $P$-module by restricting the $G$ action
to $G_0$ and decreeing that $G_+$ act trivially. Since ${\mathfrak{g}}_-$ is
Abelian, this formula agrees with (\ref{thisispartial}) for the Koszul
dif\/ferential (used in~\cite{ce} to def\/ine Lie algebra cohomology). As observed
in the appendix, this is a complex of $G_0$-modules. Alternatively, we could
come to the same conclusion by restricting attention to a standard af\/f\/ine
co\"ordinate chart $G_-\cong{\mathbb{R}}^n\hookrightarrow{\mathbb{RP}}_n$ and
noticing that the kernel of
$d:\Lambda^p\otimes{\mathbb{V}}\to\Lambda^{p+1}\otimes{\mathbb{V}}$ consists
precisely of the left-invariant ${\mathbb{V}}$-valued $p$-forms under the
action of $G_-$ on itself. Since $\nabla=d+\partial$, we may now invoke the
second realisation of $H^r({\mathfrak{g}}_-,{\mathbb{V}})$ from
Theorem~\ref{geometricrealisation} in the Appendix.

To proceed we need to be explicit concerning the irreducible representation
${\mathbb{V}}$ of ${\mathrm{SL}}(n+1,{\mathbb{R}})$. In fact, one usually deals
with complex representations of ${\mathrm{SL}}(n+1,{\mathbb{C}})$ (starting
with ${\mathfrak{sl}}(n+1,{\mathbb{C}})$) but here there is no real dif\/ference
and we shall adopt the notation from~\cite{beastwood} in denoting such
representations by attaching non-negative integers to its Dynkin diagram
\[{\mathbb{V}}=\enskip\raisebox{-5pt}{\begin{picture}(140,17)
\put(0,5){\makebox(0,0){$\bullet$}}
\put(20,5){\makebox(0,0){$\bullet$}}
\put(40,5){\makebox(0,0){$\bullet$}}
\put(60,5){\makebox(0,0){$\bullet$}}
\put(80,5){\makebox(0,0){$\bullet$}}
\put(120,5){\makebox(0,0){$\bullet$}}
\put(140,5){\makebox(0,0){$\bullet$}}
\put(0,5){\line(1,0){90}}
\put(100,5){\makebox(0,0){\scriptsize$\cdots$}}
\put(110,5){\line(1,0){30}}
\put(0,12){\makebox(0,0){\scriptsize$a_1$}}
\put(20,12){\makebox(0,0){\scriptsize$a_2$}}
\put(40,12){\makebox(0,0){\scriptsize$a_3$}}
\put(60,12){\makebox(0,0){\scriptsize$a_4$}}
\put(80,12){\makebox(0,0){\scriptsize$a_5$}}
\put(100,12){\makebox(0,0){\scriptsize$\cdots$}}
\put(120,12){\makebox(0,0){\scriptsize$a_{n-1}$}}
\put(140,12){\makebox(0,0){\scriptsize$a_n$}}
\end{picture}}\]
(meaning that $-[a_1,a_2,\dots,a_n]$ is the lowest weight of this
representation with respect to the standard basis of fundamental weights).
With this notation, here is the conclusion of Kostant's computation~\cite{k}
of Lie algebra cohomology. We describe the result as an
${\mathrm{SL}}(n,{\mathbb{R}})$-module (by attaching non-negative integers to
an $A$-series Dynkin diagram with one fewer nodes) and f\/ix the action of $G_0$
by specifying how the grading element $H\in{\mathfrak{g}}_0$ acts.
\begin{equation}\label{cohomology}
\makebox[0pt][l]{\hspace*{-20pt}$\begin{array}{rcl}
H^0({\mathfrak{g}}_-,{\mathbb{V}})&=&
\enskip\raisebox{-5pt}{\begin{picture}(120,17)
\put(0,5){\makebox(0,0){$\bullet$}}
\put(20,5){\makebox(0,0){$\bullet$}}
\put(40,5){\makebox(0,0){$\bullet$}}
\put(60,5){\makebox(0,0){$\bullet$}}
\put(100,5){\makebox(0,0){$\bullet$}}
\put(120,5){\makebox(0,0){$\bullet$}}
\put(0,5){\line(1,0){70}}
\put(80,5){\makebox(0,0){\scriptsize$\cdots$}}
\put(90,5){\line(1,0){30}}
\put(0,12){\makebox(0,0){\scriptsize$a_2$}}
\put(20,12){\makebox(0,0){\scriptsize$a_3$}}
\put(40,12){\makebox(0,0){\scriptsize$a_4$}}
\put(60,12){\makebox(0,0){\scriptsize$a_5$}}
\put(80,12){\makebox(0,0){\scriptsize$\cdots$}}
\put(100,12){\makebox(0,0){\scriptsize$a_{n-1}$}}
\put(120,12){\makebox(0,0){\scriptsize$a_n$}}
\end{picture}}\qquad
H\leadsto -c\\[5pt]
H^1({\mathfrak{g}}_-,{\mathbb{V}})&=&
\enskip\quad\raisebox{-5pt}{\begin{picture}(130,17)
\put(0,5){\makebox(0,0){$\bullet$}}
\put(30,5){\makebox(0,0){$\bullet$}}
\put(50,5){\makebox(0,0){$\bullet$}}
\put(70,5){\makebox(0,0){$\bullet$}}
\put(110,5){\makebox(0,0){$\bullet$}}
\put(130,5){\makebox(0,0){$\bullet$}}
\put(0,5){\line(1,0){80}}
\put(90,5){\makebox(0,0){\scriptsize$\cdots$}}
\put(100,5){\line(1,0){30}}
\put(0,12){\makebox(0,0){\scriptsize$a_1+a_2+1$}}
\put(30,12){\makebox(0,0){\scriptsize$a_3$}}
\put(50,12){\makebox(0,0){\scriptsize$a_4$}}
\put(70,12){\makebox(0,0){\scriptsize$a_5$}}
\put(90,12){\makebox(0,0){\scriptsize$\cdots$}}
\put(110,12){\makebox(0,0){\scriptsize$a_{n-1}$}}
\put(130,12){\makebox(0,0){\scriptsize$a_n$}}
\end{picture}}\qquad
H\leadsto -c+a_1+1\\[5pt]
H^2({\mathfrak{g}}_-,{\mathbb{V}})&=&
\enskip\raisebox{-5pt}{\begin{picture}(140,17)
\put(0,5){\makebox(0,0){$\bullet$}}
\put(30,5){\makebox(0,0){$\bullet$}}
\put(60,5){\makebox(0,0){$\bullet$}}
\put(80,5){\makebox(0,0){$\bullet$}}
\put(120,5){\makebox(0,0){$\bullet$}}
\put(140,5){\makebox(0,0){$\bullet$}}
\put(0,5){\line(1,0){90}}
\put(100,5){\makebox(0,0){\scriptsize$\cdots$}}
\put(110,5){\line(1,0){30}}
\put(0,12){\makebox(0,0){\scriptsize$a_1$}}
\put(30,12){\makebox(0,0){\scriptsize$a_2+a_3+1$}}
\put(60,12){\makebox(0,0){\scriptsize$a_4$}}
\put(80,12){\makebox(0,0){\scriptsize$a_5$}}
\put(100,12){\makebox(0,0){\scriptsize$\cdots$}}
\put(120,12){\makebox(0,0){\scriptsize$a_{n-1}$}}
\put(140,12){\makebox(0,0){\scriptsize$a_n$}}
\end{picture}}\qquad
H\leadsto -c+a_1+a_2+2\\[5pt]
H^3({\mathfrak{g}}_-,{\mathbb{V}})&=&
\enskip\raisebox{-5pt}{\begin{picture}(140,17)
\put(0,5){\makebox(0,0){$\bullet$}}
\put(20,5){\makebox(0,0){$\bullet$}}
\put(50,5){\makebox(0,0){$\bullet$}}
\put(80,5){\makebox(0,0){$\bullet$}}
\put(120,5){\makebox(0,0){$\bullet$}}
\put(140,5){\makebox(0,0){$\bullet$}}
\put(0,5){\line(1,0){90}}
\put(100,5){\makebox(0,0){\scriptsize$\cdots$}}
\put(110,5){\line(1,0){30}}
\put(0,12){\makebox(0,0){\scriptsize$a_1$}}
\put(20,12){\makebox(0,0){\scriptsize$a_2$}}
\put(50,12){\makebox(0,0){\scriptsize$a_3+a_4+1$}}
\put(80,12){\makebox(0,0){\scriptsize$a_5$}}
\put(100,12){\makebox(0,0){\scriptsize$\cdots$}}
\put(120,12){\makebox(0,0){\scriptsize$a_{n-1}$}}
\put(140,12){\makebox(0,0){\scriptsize$a_n$}}
\end{picture}}\qquad
H\leadsto -c+a_1+a_2+a_3+3\\
\vdots\vdots\hspace*{20pt}&\vdots&\hspace*{20pt}\vdots\vdots\\
H^{n-1}({\mathfrak{g}}_-,{\mathbb{V}})&=&
\enskip\raisebox{-5pt}{\begin{picture}(140,17)
\put(0,5){\makebox(0,0){$\bullet$}}
\put(20,5){\makebox(0,0){$\bullet$}}
\put(40,5){\makebox(0,0){$\bullet$}}
\put(60,5){\makebox(0,0){$\bullet$}}
\put(100,5){\makebox(0,0){$\bullet$}}
\put(140,5){\makebox(0,0){$\bullet$}}
\put(0,5){\line(1,0){70}}
\put(80,5){\makebox(0,0){\scriptsize$\cdots$}}
\put(90,5){\line(1,0){50}}
\put(0,12){\makebox(0,0){\scriptsize$a_1$}}
\put(20,12){\makebox(0,0){\scriptsize$a_2$}}
\put(40,12){\makebox(0,0){\scriptsize$a_3$}}
\put(60,12){\makebox(0,0){\scriptsize$a_4$}}
\put(80,12){\makebox(0,0){\scriptsize$\cdots$}}
\put(100,12){\makebox(0,0){\scriptsize$a_{n-2}$}}
\put(140,12){\makebox(0,0){\scriptsize$a_{n-1}+a_n+1$}}
\end{picture}}\qquad\quad
H\leadsto -c+a_1+\cdots+a_{n-1}+n-1\\[5pt]
H^n({\mathfrak{g}}_-,{\mathbb{V}})&=&
\enskip\raisebox{-5pt}{\begin{picture}(130,17)
\put(0,5){\makebox(0,0){$\bullet$}}
\put(20,5){\makebox(0,0){$\bullet$}}
\put(40,5){\makebox(0,0){$\bullet$}}
\put(60,5){\makebox(0,0){$\bullet$}}
\put(100,5){\makebox(0,0){$\bullet$}}
\put(130,5){\makebox(0,0){$\bullet$}}
\put(0,5){\line(1,0){70}}
\put(80,5){\makebox(0,0){\scriptsize$\cdots$}}
\put(90,5){\line(1,0){40}}
\put(0,12){\makebox(0,0){\scriptsize$a_1$}}
\put(20,12){\makebox(0,0){\scriptsize$a_2$}}
\put(40,12){\makebox(0,0){\scriptsize$a_3$}}
\put(60,12){\makebox(0,0){\scriptsize$a_4$}}
\put(80,12){\makebox(0,0){\scriptsize$\cdots$}}
\put(100,12){\makebox(0,0){\scriptsize$a_{n-2}$}}
\put(130,12){\makebox(0,0){\scriptsize$a_{n-1}$}}
\end{picture}}\qquad
H\leadsto -c+a_1+\cdots+a_{n-1}+a_n+n
\end{array}$}\end{equation}
where
\[c=\frac{na_1+(n-1)a_2+(n-2)a_3+\cdots+2a_{n-1}+a_n}{n+1}\]
(obtained by acting on $[a_1,a_2,\dots,a_n]$ with the f\/irst column of the
inverse Cartan matrix for ${\mathfrak{sl}}(n+1)$). Notice that each cohomology
is an irreducible representation of~$G_0$.

\begin{theorem}\label{E1level}
The $E_1$-level of the spectral sequence of the filtered complex
$\nabla:\Lambda^\bullet\otimes V$ consists of irreducible homogeneous vector
bundles on~${\mathbb{RP}}_n$ under the action of
$G={\mathrm{SL}}(n+1,{\mathbb{R}})$. Only the following terms are non-zero
\[\begin{array}{@{}rcl}
E_1^{0,0}&\leftrightsquigarrow&H^0({\mathfrak{g}}_-,{\mathbb{V}})\\
E_1^{a_1+1,-a_1}&\leftrightsquigarrow&H^1({\mathfrak{g}}_-,{\mathbb{V}})\\
E_1^{a_1+a_2+2,-a_1-a_2}&\leftrightsquigarrow&
H^2({\mathfrak{g}}_-,{\mathbb{V}})\\
E_1^{a_1+a_2+a_3+3,-a_1-a_2-a_3}&\leftrightsquigarrow&
H^3({\mathfrak{g}}_-,{\mathbb{V}})\\
\vdots\vdots\hspace*{20pt}&\vdots&\hspace*{20pt}\vdots\vdots\\
E_1^{N+n-1-a_n,-N+a_n}\enskip=\enskip
E_1^{a_1+a_2+a_3+\cdots+a_{n-1}+n-1,-a_1-a_2-a_3-\cdots-a_{n-1}}
&\leftrightsquigarrow&
H^{n-1}({\mathfrak{g}}_-,{\mathbb{V}})\\
E_1^{N+n,-N}\enskip=\enskip
E_1^{a_1+a_2+a_3+\cdots+a_{n-1}+a_n+n,-a_1-a_2-a_3-\cdots-a_{n-1}-a_n}
&\leftrightsquigarrow&
H^n({\mathfrak{g}}_-,{\mathbb{V}})
\end{array}\]
meaning that the bundle in question is induced by the $G_0$-module as listed
and extended trivially as a $P$-module.
\end{theorem}

\begin{proof} According to (\ref{E0pq}), we already know that
$\Lambda^r\otimes V$ is spread along the
$r^{\mathrm{th}}$ diagonal
\[\Lambda^r\otimes V=\Lambda^r\otimes V_0+\Lambda^r\otimes
V_1+\Lambda^r\otimes V_2+\cdots =
E_0^{r,0}+E_0^{r+1,-1}+E_0^{r+2,-2}+\cdots\]
of the $E_0$-level of the spectral sequence and that the $E_0$-dif\/ferential is
induced by the Koszul dif\/ferential~(\ref{E0differential}) def\/ining the Lie
algebra cohomology $H^r({\mathfrak{g}}_-,{\mathbb{V}})$. We see from
(\ref{cohomology}) that each of these $H^r({\mathfrak{g}}_-,{\mathbb{V}})$ is
irreducible and so it follows that the $r^{\mathrm{th}}$ diagonal of the
$E_1$-level consists of a single irreducible homogeneous vector bundle. To
complete the proof it suf\/f\/ices to locate the position of this bundle along this
particular diagonal and, to do this, the action of the grading element $H$
turns out to be suf\/f\/icient. More precisely,
\[\begin{array}{@{}rclrcl}
H\mbox{ acts by}&-c&\mbox{on }V_0&
\therefore\;H\mbox{ acts by}&-c+k&\mbox{on }V_k,\\
H\mbox{ acts by}&-1&\mbox{on }{\mathfrak{g}}_-&
\therefore\;H\mbox{ acts by}&p&\mbox{on }\Lambda^p{\mathfrak{g}}_-^*,\\
&&&\therefore\;H\mbox{ acts by}&-c+k+p&\mbox{on }
\Lambda^p{\mathfrak{g}}_-^*\otimes{\mathbb{V}}_k,\\
&&&\therefore\;H\mbox{ acts by}&-c+p&\mbox{on }
\Lambda^{p+q}{\mathfrak{g}}_-^*\otimes{\mathbb{V}}_{-q}\leftrightsquigarrow
E_0^{p,q}.
\end{array}\]
Therefore $H$ acts by $-c+p$ on the $G_0$-module inducing $E_1^{p,q}$ and so,
from the action of $H$ in table~(\ref{cohomology}), the bundle induced by
$H^r({\mathfrak{g}}_-,{\mathbb{V}})$ is located at
$E_1^{a_1+a_2+\cdots+a_r+r,-a_1-a_2-\cdots-a_r}$.

Finally, we are claiming that $N=a_1+a_2+\cdots+a_n$. To see this we note that,
since $H$ acts on $V_0$ by $-c$, it acts on $V_N$ by $N-c$. On the other hand,
in accordance with the action of the longest element of the Weyl group,
\[{\mathbb{V}}^*=\enskip\raisebox{-5pt}{\begin{picture}(140,17)
\put(0,5){\makebox(0,0){$\bullet$}}
\put(20,5){\makebox(0,0){$\bullet$}}
\put(60,5){\makebox(0,0){$\bullet$}}
\put(80,5){\makebox(0,0){$\bullet$}}
\put(100,5){\makebox(0,0){$\bullet$}}
\put(120,5){\makebox(0,0){$\bullet$}}
\put(140,5){\makebox(0,0){$\bullet$}}
\put(0,5){\line(1,0){30}}
\put(40,5){\makebox(0,0){\scriptsize$\cdots$}}
\put(50,5){\line(1,0){90}}
\put(0,12){\makebox(0,0){\scriptsize$a_n$}}
\put(20,12){\makebox(0,0){\scriptsize$a_{n-1}$}}
\put(40,12){\makebox(0,0){\scriptsize$\cdots$}}
\put(60,12){\makebox(0,0){\scriptsize$a_5$}}
\put(80,12){\makebox(0,0){\scriptsize$a_4$}}
\put(100,12){\makebox(0,0){\scriptsize$a_3$}}
\put(120,12){\makebox(0,0){\scriptsize$a_2$}}
\put(140,12){\makebox(0,0){\scriptsize$a_1$}}
\end{picture}}\enskip={\mathbb{V}}_N^*+{\mathbb{V}}_{N-1}^*+\cdots+
{\mathbb{V}}_2^*+{\mathbb{V}}_1^*+{\mathbb{V}}_0^*\]
and so $H$ acts on ${\mathbb{V}}_N^*$ by $-c^\prime$, where
\[c^\prime=\frac{na_n+(n-1)a_{n-1}+\cdots+3a_3+2a_2+a_1}{n+1}.\]
We conclude that $N=c+c^\prime=a_1+a_2+a_3+\cdots+a_{n-1}+a_n$, as required.
\end{proof}

As outlined in Section~\ref{outline}, the $E_1$-level of this spectral sequence is
rather sparse~(\ref{sparse}) and Theorem~\ref{E1level} says exactly how sparse.
In particular, since there is only one non-zero bundle on each diagonal of the
$E_1$-level, the general theory of spectral sequences provides a complex of
dif\/ferential operators
\[H^0\to H^1\to H^2\to H^3\to\cdots\to H^{n-1}\to H^n\to 0\]
whose cohomology on the level of sheaves coincides with that of
(\ref{coupled}), namely ${\mathbb{V}}$ for $\ker:H^0\to H^1$ and otherwise
zero. The bundle $H^r$ is induced on ${\mathbb{RP}}_n=G/P$ from the
$G_0$-module $H^r({\mathfrak{g}}_-,{\mathbb{V}})$ extended trivially as as
$P$-module. This is the BGG resolution (constructed on a general $G/P$ by
Lepowsky~\cite{l} on the level of generalised Verma modules). More explicitly,
in the discussion following Corollary~\ref{easypeasy}, in any af\/f\/ine
co\"ordinate patch we identif\/ied
$\nabla:\Lambda^p\otimes V\to\Lambda^{p+1}\otimes V$ as $d+\partial$ where
$V|_{{\mathbb{R}}^n}$ is trivialised as ${\mathbb{R}}^n\times{\mathbb{V}}$ and
\[E_0^{p,q}=\Lambda^{p+q}\otimes V_{-q}\cong
\Lambda^{p+q}\otimes{\mathbb{V}}_{-q}\stackrel{d}{\longrightarrow}
\Lambda^{p+q+1}\otimes{\mathbb{V}}_{-q}\cong\Lambda^{p+q+1}\otimes V_{-q}=
E_0^{p+1,q}\]
is the exterior derivative with values in~${\mathbb{V}}_{-q}$. Adding these
operators to the diagram (\ref{E0}) gives a~double complex as employed by
Baston~\cite[p.~120]{b}. We conclude that in any af\/f\/ine co\"ordinate patch, our
spectral sequence (of a f\/iltered complex) coincides with Baston's spectral
sequence (of a double complex). In particular, the operators $H^k\to H^{k+1}$
are obtained as zigzag compositions of the f\/irst order dif\/ferential operators
$d$ together with choices of algebraic splittings of the operators~$\partial$
(Baston and others use the algebraic adjoint $\partial^*$ introduced by
Kostant~\cite{k}). This conf\/irms that the resulting operators $H^k\to H^{k+1}$
are dif\/ferential. By using the spectral sequence of a f\/iltered complex as we
have done, it is manifest that the dif\/ferential operators $H^k\to H^{k+1}$ are
independent of any choice of splittings and that the whole construction and
resulting BGG complex is $G$-equivariant. It is also clear that the f\/irst steps
in this approach can be taken on an arbitrary homogeneous space~$G/P$.

\section{Formul{\ae} for the BGG operators}
Already, the af\/f\/ine invariance of the operators occurring in the BGG resolution
on ${\mathbb{RP}}_n$ f\/ixes their formul{\ae} with respect to the f\/lat
connection on ${\mathbb{R}}^n\subset{\mathbb{RP}}_n$ as follows (af\/f\/ine
invariance being ensured by noting that $Q$ as in (\ref{Q}) is acting on
${\mathbb{R}}^n=Q/G_0$ by af\/f\/ine transformations). Firstly, consider the symbol
of the BGG dif\/ferential operator $H^r\to H^{r+1}$. As a homomorphism of
homogeneous bundles $\bigodot^s\!\Lambda^1\otimes H^r\to H^{r+1}$ for
some~$s$, it is induced by a homomorphism of $G_0$-modules
\[\textstyle
\bigodot^s\!{\mathfrak{g}}_-^*\otimes H^r({\mathfrak{g}}_-,{\mathbb{V}})\to
H^{r+1}({\mathfrak{g}}_-,{\mathbb{V}})\]
for some $s$, which we can determine just from the action of the grading
element. Specif\/ically, we know that
\[\begin{array}{@{}rcl}
H\mbox{ acts by}&-c+a_1+a_2+\cdots+a_r+r&
\mbox{on }H^r({\mathfrak{g}}_-,{\mathbb{V}})\\
H\mbox{ acts by}&-1&
\mbox{on }{\mathfrak{g}}_-\\
\therefore\enskip H\mbox{ acts by}&s&
\mbox{on }\bigodot^s\!{\mathfrak{g}}_-^*.
\end{array}\]
It is immediate that $s=a_{r+1}+1$. Furthermore, from the
Littlewood--Richardson rules~(e.g.,~\cite{fh}), as an
${\mathfrak{sl}}(n,{\mathbb{R}})$-module
$H^{r+1}({\mathfrak{g}}_-,{\mathbb{V}})$ occurs with multiplicity one in the
following decomposition
\[\begin{array}{@{}lcc}
\hspace*{21.4pt}\bigodot^{a_{r+1}+1}\!{\mathfrak{g}}_-^*
\otimes H^r({\mathfrak{g}}_-,{\mathbb{V}})
&=&\cdots\oplus H^{r+1}({\mathfrak{g}}_-,{\mathbb{V}})\oplus\cdots\\
\hspace*{79.5pt}\|&&\|\\
\raisebox{-5pt}{\begin{picture}(70,17)
\put(0,5){\makebox(0,0){$\bullet$}}
\put(30,5){\makebox(0,0){$\bullet$}}
\put(70,5){\makebox(0,0){$\bullet$}}
\put(0,5){\line(1,0){40}}
\put(50,5){\makebox(0,0){\scriptsize$\cdots$}}
\put(60,5){\line(1,0){10}}
\put(0,12){\makebox(0,0){\scriptsize$a_{r+1}+1$}}
\put(30,12){\makebox(0,0){\scriptsize$0$}}
\put(50,12){\makebox(0,0){\scriptsize$\cdots$}}
\put(70,12){\makebox(0,0){\scriptsize$0$}}
\end{picture}}\enskip\otimes\enskip
\raisebox{-5pt}{\begin{picture}(140,17)
\put(0,5){\makebox(0,0){$\bullet$}}
\put(60,5){\makebox(0,0){$\bullet$}}
\put(100,5){\makebox(0,0){$\bullet$}}
\put(140,5){\makebox(0,0){$\bullet$}}
\put(0,5){\line(1,0){10}}
\put(20,5){\makebox(0,0){\scriptsize$\cdots$}}
\put(30,5){\line(1,0){80}}
\put(120,5){\makebox(0,0){\scriptsize$\cdots$}}
\put(130,5){\line(1,0){10}}
\put(0,12){\makebox(0,0){\scriptsize$a_1$}}
\put(20,12){\makebox(0,0){\scriptsize$\cdots$}}
\put(60,12){\makebox(0,0){\scriptsize$a_r+a_{r+1}+1$}}
\put(100,12){\makebox(0,0){\scriptsize$a_{r+2}$}}
\put(120,12){\makebox(0,0){\scriptsize$\cdots$}}
\put(140,12){\makebox(0,0){\scriptsize$a_n$}}
\end{picture}}&=&
\oplus\enskip\raisebox{-5pt}{\begin{picture}(140,17)
\put(0,5){\makebox(0,0){$\bullet$}}
\put(40,5){\makebox(0,0){$\bullet$}}
\put(80,5){\makebox(0,0){$\bullet$}}
\put(140,5){\makebox(0,0){$\bullet$}}
\put(0,5){\line(1,0){10}}
\put(20,5){\makebox(0,0){\scriptsize$\cdots$}}
\put(30,5){\line(1,0){80}}
\put(120,5){\makebox(0,0){\scriptsize$\cdots$}}
\put(130,5){\line(1,0){10}}
\put(0,12){\makebox(0,0){\scriptsize$a_1$}}
\put(20,12){\makebox(0,0){\scriptsize$\cdots$}}
\put(40,12){\makebox(0,0){\scriptsize$a_r$}}
\put(80,12){\makebox(0,0){\scriptsize$a_{r+1}+a_{r+2}+1$}}
\put(120,12){\makebox(0,0){\scriptsize$\cdots$}}
\put(140,12){\makebox(0,0){\scriptsize$a_n$}}
\end{picture}}\enskip\oplus
\end{array}\]
The symbol of the BGG operator $H^r\to H^{r+1}$
is thus determined uniquely (up to scale) as induced by the projection onto
this summand. Similarly, there are no lower order terms since no appropriate
invariant homomorphisms are available. We record this conclusion as follows.
\begin{theorem}
Each bundle $H^r$ in the BGG complex on ${\mathbb{RP}}_n$ is an irreducible
tensor bundle and the differential operator $\nabla^{(s)}:H^r\to H^{r+1}$ is
given by $\phi\mapsto\pi(\nabla^s\phi)$ where $\nabla$ is the flat affine
connection and $\pi:\bigodot^s\!\Lambda^1\otimes H^r\to H^{r+1}$ is the unique
projection onto this irreducible summand.
\end{theorem}
Equivalently, the bundles $H^r$ and the operators between them are exhibited as
Young tableau in~\cite{eProjective}. It is often useful, however, to be
able to write the BGG complex globally on ${\mathbb{RP}}_n$ without recourse to
af\/f\/ine co\"ordinates and projective invariance. For this, we shall use the
round metric on the sphere and the corresponding Levi Civita connection on
$S^n$ or~${\mathbb{RP}}_n$. By way of normalisation, if $g_{ab}$ denotes the
round metric and $\delta_a{}^b$ the identity endomorphism on the tangent
bundle, let us choose the radius of the sphere so that
\begin{equation}\label{cc}
R_{ab}{}^c{}_d=\delta_a{}^cg_{bd}-\delta_b{}^cg_{ad}\qquad\mbox{where}\quad
(\nabla_a\nabla_b-\nabla_b\nabla_a)X^d=R_{ab}{}^c{}_dX^d\end{equation}
def\/ines the Riemann curvature tensor. For any irreducible covariant tensor
bundle $E$ on ${\mathbb{RP}}_n$, let us write $\nabla^2$ for the composition
\[\textstyle\bigodot^r\!\Lambda^1\otimes E\xrightarrow{\,\nabla\circ\nabla\,}
\Lambda^1\otimes\Lambda^1\otimes\bigodot^r\!\Lambda^1\otimes E\to
\bigodot^{r+2}\!\Lambda^1\otimes E\]
and $g$ for the composition
\[\textstyle\bigodot^r\!\Lambda^1\otimes E
\xrightarrow{\,g\otimes{\mathrm{Id}}\,}
\bigodot^2\!\Lambda^1\otimes\bigodot^r\!\Lambda^1\otimes E\to
\bigodot^{r+2}\!\Lambda^1\otimes E.\]
\begin{theorem}\label{round_formulae}
If $s$ is odd, the operator $\nabla^{(s)}:H^r\to H^{r+1}$ is given by
\[\pi\big((\nabla^2+(s-1)^2g)\cdots(\nabla^2+16g)(\nabla^2+4g)\nabla\big)\]
and, if $s$ is even,
\[\pi\big((\nabla^2+(s-1)^2g)\cdots(\nabla^2+9g)(\nabla^2+g)\big).\]
\end{theorem}
\begin{proof}
It is convenient to use some notation from~\cite{eProjective} where the
projective BGG operators are written as acting between covariant tensor bundles
specif\/ied by weighted Young tableau: %--
\[\begin{picture}(170,35)(0,85)
\put(0,115){\line(1,0){150}}
\put(0,110){\line(1,0){150}}
\put(0,105){\line(1,0){120}}
\put(0,100){\line(1,0){90}}
\put(0,95){\line(1,0){60}}
\put(0,90){\line(1,0){30}}
\put(0,90){\line(0,1){25}}
\put(30,90){\line(0,1){25}}
\put(60,95){\line(0,1){20}}
\put(90,100){\line(0,1){15}}
\put(120,105){\line(0,1){10}}
\put(150,110){\line(0,1){5}}
\put(15,92.5){\makebox(0,0){$\cdots$}}
\put(15,97.5){\makebox(0,0){$\cdots$}}
\put(15,102.5){\makebox(0,0){$\cdots$}}
\put(15,107.5){\makebox(0,0){$\cdots$}}
\put(15,112.5){\makebox(0,0){$\cdots$}}
\put(45,97.5){\makebox(0,0){$\cdots$}}
\put(45,102.5){\makebox(0,0){$\cdots$}}
\put(45,107.5){\makebox(0,0){$\cdots$}}
\put(45,112.5){\makebox(0,0){$\cdots$}}
\put(75,102.5){\makebox(0,0){$\cdots$}}
\put(75,107.5){\makebox(0,0){$\cdots$}}
\put(75,112.5){\makebox(0,0){$\cdots$}}
\put(105,107.5){\makebox(0,0){$\cdots$}}
\put(105,112.5){\makebox(0,0){$\cdots$}}
\put(135,112.5){\makebox(0,0){$\cdots$}}
\put(5,90){\line(0,1){25}}
\put(25,90){\line(0,1){25}}
\put(35,95){\line(0,1){20}}
\put(55,95){\line(0,1){20}}
\put(65,100){\line(0,1){15}}
\put(85,100){\line(0,1){15}}
\put(95,105){\line(0,1){10}}
\put(115,105){\line(0,1){10}}
\put(125,110){\line(0,1){5}}
\put(145,110){\line(0,1){5}}
\put(40,85){\vector(-1,0){40}}
\put(50,85){\vector(1,0){40}}
\put(45,85){\makebox(0,0){$\scriptstyle b$}}
\put(155,100){\makebox(0,0){$\scriptstyle (w)$}}
\end{picture}
\raisebox{15pt}{$\stackrel{\nabla^{(s)}}{\verylongrightarrow}\quad$}
\begin{picture}(170,35)(0,85)
\put(0,115){\line(1,0){150}}
\put(0,110){\line(1,0){150}}
\put(0,105){\line(1,0){120}}
\put(0,100){\line(1,0){90}}
\put(0,95){\line(1,0){60}}
\put(0,90){\line(1,0){30}}
\put(0,90){\line(0,1){25}}
\put(30,90){\line(0,1){25}}
\put(60,95){\line(0,1){20}}
\put(90,100){\line(0,1){15}}
\put(120,105){\line(0,1){10}}
\put(150,110){\line(0,1){5}}
\put(15,92.5){\makebox(0,0){$\cdots$}}
\put(15,97.5){\makebox(0,0){$\cdots$}}
\put(15,102.5){\makebox(0,0){$\cdots$}}
\put(15,107.5){\makebox(0,0){$\cdots$}}
\put(15,112.5){\makebox(0,0){$\cdots$}}
\put(45,97.5){\makebox(0,0){$\cdots$}}
\put(45,102.5){\makebox(0,0){$\cdots$}}
\put(45,107.5){\makebox(0,0){$\cdots$}}
\put(45,112.5){\makebox(0,0){$\cdots$}}
\put(75,102.5){\makebox(0,0){$\cdots$}}
\put(75,107.5){\makebox(0,0){$\cdots$}}
\put(75,112.5){\makebox(0,0){$\cdots$}}
\put(105,107.5){\makebox(0,0){$\cdots$}}
\put(105,112.5){\makebox(0,0){$\cdots$}}
\put(135,112.5){\makebox(0,0){$\cdots$}}
\put(5,90){\line(0,1){25}}
\put(25,90){\line(0,1){25}}
\put(35,95){\line(0,1){20}}
\put(55,95){\line(0,1){20}}
\put(65,100){\line(0,1){15}}
\put(85,100){\line(0,1){15}}
\put(95,105){\line(0,1){10}}
\put(115,105){\line(0,1){10}}
\put(125,110){\line(0,1){5}}
\put(145,110){\line(0,1){5}}
\put(40,85){\vector(-1,0){40}}
\put(60,85){\vector(1,0){52}}
\put(50,85){\makebox(0,0){$\scriptstyle b+s$}}
\put(155,100){\makebox(0,0){$\scriptstyle (w)\makebox[0pt][l]{$\,.$}$}}
\thicklines
\put(90,100){\line(1,0){22}}
\put(90,105){\line(1,0){22}}
\put(90,100){\line(0,1){5}}
\put(95,100){\line(0,1){5}}
\put(107,100){\line(0,1){5}}
\put(112,100){\line(0,1){5}}
\put(98,102.5){\makebox(0,0){.}}
\put(101,102.5){\makebox(0,0){.}}
\put(104,102.5){\makebox(0,0){.}}
\end{picture}\]
Here, if it is the $r^{\mathrm{th}}$ row to which the boxes on the right hand
side are being added and there are a total of $v$ boxes on the left hand side
(so that $v$ is the valence of the corresponding tensor), then $w+r=s+v+b$
(see~\cite{eProjective}). In
particular, if we consider only f\/irst order operators
\begin{equation}\label{firstorder}\begin{picture}(170,35)(0,85)
\put(0,115){\line(1,0){150}}
\put(0,110){\line(1,0){150}}
\put(0,105){\line(1,0){120}}
\put(0,100){\line(1,0){90}}
\put(0,95){\line(1,0){60}}
\put(0,90){\line(1,0){30}}
\put(0,90){\line(0,1){25}}
\put(30,90){\line(0,1){25}}
\put(60,95){\line(0,1){20}}
\put(90,100){\line(0,1){15}}
\put(120,105){\line(0,1){10}}
\put(150,110){\line(0,1){5}}
\put(15,92.5){\makebox(0,0){$\cdots$}}
\put(15,97.5){\makebox(0,0){$\cdots$}}
\put(15,102.5){\makebox(0,0){$\cdots$}}
\put(15,107.5){\makebox(0,0){$\cdots$}}
\put(15,112.5){\makebox(0,0){$\cdots$}}
\put(45,97.5){\makebox(0,0){$\cdots$}}
\put(45,102.5){\makebox(0,0){$\cdots$}}
\put(45,107.5){\makebox(0,0){$\cdots$}}
\put(45,112.5){\makebox(0,0){$\cdots$}}
\put(75,102.5){\makebox(0,0){$\cdots$}}
\put(75,107.5){\makebox(0,0){$\cdots$}}
\put(75,112.5){\makebox(0,0){$\cdots$}}
\put(105,107.5){\makebox(0,0){$\cdots$}}
\put(105,112.5){\makebox(0,0){$\cdots$}}
\put(135,112.5){\makebox(0,0){$\cdots$}}
\put(5,90){\line(0,1){25}}
\put(25,90){\line(0,1){25}}
\put(35,95){\line(0,1){20}}
\put(55,95){\line(0,1){20}}
\put(65,100){\line(0,1){15}}
\put(85,100){\line(0,1){15}}
\put(95,105){\line(0,1){10}}
\put(115,105){\line(0,1){10}}
\put(125,110){\line(0,1){5}}
\put(145,110){\line(0,1){5}}
\put(40,85){\vector(-1,0){40}}
\put(50,85){\vector(1,0){40}}
\put(45,85){\makebox(0,0){$\scriptstyle b$}}
\put(155,100){\makebox(0,0){$\scriptstyle (w)$}}
\end{picture}
\raisebox{15pt}{$\stackrel{\nabla}{\verylongrightarrow}\quad$}
\begin{picture}(170,35)(0,85)
\put(0,115){\line(1,0){150}}
\put(0,110){\line(1,0){150}}
\put(0,105){\line(1,0){120}}
\put(0,100){\line(1,0){90}}
\put(0,95){\line(1,0){60}}
\put(0,90){\line(1,0){30}}
\put(0,90){\line(0,1){25}}
\put(30,90){\line(0,1){25}}
\put(60,95){\line(0,1){20}}
\put(90,100){\line(0,1){15}}
\put(120,105){\line(0,1){10}}
\put(150,110){\line(0,1){5}}
\put(15,92.5){\makebox(0,0){$\cdots$}}
\put(15,97.5){\makebox(0,0){$\cdots$}}
\put(15,102.5){\makebox(0,0){$\cdots$}}
\put(15,107.5){\makebox(0,0){$\cdots$}}
\put(15,112.5){\makebox(0,0){$\cdots$}}
\put(45,97.5){\makebox(0,0){$\cdots$}}
\put(45,102.5){\makebox(0,0){$\cdots$}}
\put(45,107.5){\makebox(0,0){$\cdots$}}
\put(45,112.5){\makebox(0,0){$\cdots$}}
\put(75,102.5){\makebox(0,0){$\cdots$}}
\put(75,107.5){\makebox(0,0){$\cdots$}}
\put(75,112.5){\makebox(0,0){$\cdots$}}
\put(105,107.5){\makebox(0,0){$\cdots$}}
\put(105,112.5){\makebox(0,0){$\cdots$}}
\put(135,112.5){\makebox(0,0){$\cdots$}}
\put(5,90){\line(0,1){25}}
\put(25,90){\line(0,1){25}}
\put(35,95){\line(0,1){20}}
\put(55,95){\line(0,1){20}}
\put(65,100){\line(0,1){15}}
\put(85,100){\line(0,1){15}}
\put(95,105){\line(0,1){10}}
\put(115,105){\line(0,1){10}}
\put(125,110){\line(0,1){5}}
\put(145,110){\line(0,1){5}}
\put(40,85){\vector(-1,0){40}}
\put(60,85){\vector(1,0){35}}
\put(50,85){\makebox(0,0){$\scriptstyle b+1$}}
\put(155,100){\makebox(0,0){$\scriptstyle (w)\makebox[0pt][l]{$\,,$}$}}
\thicklines
\put(90,100){\line(1,0){5}}
\put(90,105){\line(1,0){5}}
\put(90,100){\line(0,1){5}}
\put(95,100){\line(0,1){5}}
\end{picture}\end{equation}
then $w=v+b+1-r$. More generally, taking conventions from
\cite{eProjective}, if $\nabla$ and
$\hat\nabla$ are two torsion-free connections in the same projective class
\[\hat\nabla_a\phi_b=\nabla_a\phi_b-\Upsilon_a\phi_b-\Upsilon_b\phi_a\]
for some $1$-form~$\Upsilon_a$, then for $\phi_{bc\cdots d}$ having symmetries
and projective weight specif\/ied by the left hand side of~(\ref{firstorder}),
\[\pi(\hat\nabla_a\phi_{bc\cdots d})=
\pi\big(\nabla_a\phi_{bc\cdots d}
+(w-(v+b+1-r))\Upsilon_a\phi_{bc\cdots d}\big),\]
where $\pi$ is the Young projector corresponding to the right hand side
of~(\ref{firstorder}). We may iterate this formula, adding more boxes to the
$r^{\mathrm{th}}$ row (assuming that there is room to do so) and, each time,
both $b$ and $v$ increase by one. Suppressing indices, after $s$
iterations the result is that
\[\pi(\hat\nabla^s\phi)=\pi\big(
(\nabla+(k-2s+2)\Upsilon\phi)\cdots
(\nabla+(k-4)\Upsilon\phi)(\nabla+(k-2)\Upsilon\phi)
(\nabla+k\Upsilon)\phi\big),\]
where we are writing $k=w-(v+b+1-r)$. In particular, if $s=w+r-v-b$ as it
is in the case of a BGG operator, then this iteration reads
\[\pi(\hat\nabla^s\phi)=\pi\big(
(\nabla-k\Upsilon\phi)\cdots
(\nabla+(k-4)\Upsilon\phi)(\nabla+(k-2)\Upsilon\phi)
(\nabla+k\Upsilon)\phi\big),\]
where $k=s-1$.
So far, this conclusion holds under any projective change of connection but now
we specialise to the case of the round connection $\nabla$ on the sphere, being
projectively equivalent to the f\/lat connection $\hat\nabla$ (under gnomonic
projection). The general formula \cite[(3.4)]{eProjective} for the
change in the Ricci tensor specialises to
\[0=g_{ab}-\nabla_a\Upsilon_b+\Upsilon_a\Upsilon_b\qquad\mbox{or, suppressing
indices,}\qquad\nabla\Upsilon=g+\Upsilon^2.\]
In this equation $\Upsilon$ is viewed as a tensor but if it is viewed as an
operator $\phi\stackrel{\Upsilon}{\longmapsto}\Upsilon\phi$, then we should
write $\nabla\Upsilon=\Upsilon\nabla+g+\Upsilon^2$. This equation allows us to
deal with the iterated formula above. For example, when $k=2$, also bearing in
mind that as operators $\nabla$ and $g$ commute,
 \begin{gather*}
 \underline{\pi\big(\hat\nabla^3\phi\big)}=
 \pi\big((\nabla-2\Upsilon)\nabla(\nabla+2\Upsilon)\big)\\
\hphantom{\underline{\pi\big(\hat\nabla^3\phi\big)}}{} = \pi\big(\nabla^3-2\Upsilon\nabla^2+2\nabla(\nabla\Upsilon)
-4\Upsilon(\nabla\Upsilon)\big)\\
\hphantom{\underline{\pi\big(\hat\nabla^3\phi\big)}}{} =
\pi\big(\nabla^3-2\Upsilon\nabla^2+2\nabla\big(\Upsilon\nabla+g+\Upsilon^2\big)
-4\Upsilon\big(\Upsilon\nabla+g+\Upsilon^2\big)\big)\\
\hphantom{\underline{\pi\big(\hat\nabla^3\phi\big)}}{} = \pi\big(\nabla^3-2\Upsilon\nabla^2+2(\nabla\Upsilon)\nabla+2g\nabla
+2(\nabla\Upsilon)\Upsilon
-4\Upsilon^2\nabla-4g\Upsilon-4\Upsilon^3\big)\\
\hphantom{\underline{\pi\big(\hat\nabla^3\phi\big)}}{} = \pi\big(\nabla^3-2\Upsilon\nabla^2+2\big(\Upsilon\nabla+g+\Upsilon^2\big)\nabla+2g\nabla
+2(\nabla\Upsilon)\Upsilon
-4\Upsilon^2\nabla-4g\Upsilon-4\Upsilon^3\big)\\
\hphantom{\underline{\pi\big(\hat\nabla^3\phi\big)}}{} = \pi\big(\nabla^3+4g\nabla+2\big(\Upsilon\nabla+g+\Upsilon^2\big)\Upsilon
-2\Upsilon^2\nabla-4g\Upsilon-4\Upsilon^3\big)\\
\hphantom{\underline{\pi\big(\hat\nabla^3\phi\big)}}{} = \pi\big(\nabla^3+4g\nabla+2\Upsilon(\nabla\Upsilon)
-2\Upsilon^2\nabla-2g\Upsilon-2\Upsilon^3\big)\\
\hphantom{\underline{\pi\big(\hat\nabla^3\phi\big)}}{} = \pi\big(\nabla^3+4g\nabla+2\Upsilon\big(\Upsilon\nabla+g+\Upsilon^2\big)
-2\Upsilon^2\nabla-2g\Upsilon-2\Upsilon^3\big)\\
\hphantom{\underline{\pi\big(\hat\nabla^3\phi\big)}}{} = \pi(\nabla^3+4g\nabla) =
\underline{\pi\big(\big(\nabla^2+4g\big)\nabla\big)},
\end{gather*}
as advertised in the statement of the theorem. For higher~$k$, direct
calculations rapidly get out of hand. Instead, it suf\/f\/ices to prove the
following lemma in which we have isolated the required algebra (and then we
prove the lemma by indirect means).
\end{proof}

\begin{remark}Our curvature normalisation (\ref{cc}) implies that the Ricci
tensor on our round sphere is given by
\[R_{ab}\equiv R_{ca}{}^c{}_b=(n-1)g_{ab}.\]
Thus, the metric $g_{ab}$ coincides with $\frac{1}{n-1}R_{ab}$, a tensor
generally known in projective dif\/ferential geometry \cite{eProjective} as the
{\em Schouten tensor\/} or {\em Rho-tensor\/}~$\Rho_{ab}$. Replacing $g$ by
$\Rho$ in the formul{\ae} of Theorem~\ref{round_formulae} gives expressions
that are valid on any space of constant curvature. More generally, there is a
Rho-tensor that arises in similar contexts~\cite{cds} within parabolic
dif\/ferential geometry~\cite{thebook}.
\end{remark}
\begin{lemma}\label{pretricky} Define an associative algebra
${\mathcal{R}}={\mathbb{R}}\langle\nabla,\Upsilon\rangle$ with generators
subject to the `Riccati relation'
$\nabla\Upsilon=\Upsilon\nabla+1+\Upsilon^2$. Then the following identities
hold in~${\mathcal{R}}$. If $k$ is even, then
\begin{gather*}
(\nabla-k\Upsilon)(\nabla-(k-2)\Upsilon)\cdots
(\nabla+(k-2)\Upsilon)(\nabla+k\Upsilon)=
(\nabla^2+k^2)\cdots(\nabla^2+16)(\nabla^2+4)\nabla\!
\end{gather*}
and, if $k$ is odd, then
\begin{gather*}
(\nabla-k\Upsilon)(\nabla-(k-2)\Upsilon)\cdots
(\nabla+(k-2)\Upsilon)(\nabla+k\Upsilon)=
(\nabla^2+k^2)\cdots(\nabla^2+9)(\nabla^2+1).
\end{gather*}
\end{lemma}
\begin{proof} The algebra ${\mathcal{R}}$ may be realised by the
following dif\/ferential operators on the circle
\[f(\theta)\stackrel{\nabla}{\longmapsto} df(\theta)/d\theta,\qquad
f(\theta)\stackrel{\Upsilon}{\longmapsto}(\tan\theta)f(\theta).\]
To see this, note that these operators certainly satisfy the Riccati
relation and we are required, therefore, to show that they satisfy no further
relations. Within ${\mathcal{R}}$ we may normalise any element as follows. By
induction, the Riccati relation extends to
\[\nabla\Upsilon^\ell=
\Upsilon^\ell\nabla+\ell\big(\Upsilon^{\ell-1}+\Upsilon^{\ell+1}\big),\qquad
\forall\,\ell\geq 1\]
whence
\[\nabla^k\Upsilon^\ell
=\nabla^{k-1}\big(\Upsilon^\ell\nabla
+\ell\big(\Upsilon^{\ell-1}+\Upsilon^{\ell+1}\big)\big)
=\big(\nabla^{k-1}\Upsilon^\ell\big)\nabla
+\ell\big(\nabla^{k-1}\Upsilon^{\ell-1}\big)+\ell\big(\nabla^{k-1}\Upsilon^{\ell+1}\big)\]
and it follows by induction on $k$ that
\begin{gather*}
\nabla^k\Upsilon^\ell=
\Upsilon^\ell\nabla^k+k\ell\big(\Upsilon^{\ell-1}+\Upsilon^{\ell+1}\big)\nabla^{k-1}\\
\hphantom{\nabla^k\Upsilon^\ell=}{}
+\tfrac12k(k-1)\ell\big((\ell-1)\Upsilon^{\ell-2}+2\ell\Upsilon^\ell
+(\ell+1)\Upsilon^{\ell+2}\big)\nabla^{k-2}+\cdots,
\end{gather*}
where the ellipsis $\cdots$ denotes terms of lower order in $\nabla$ with
coef\/f\/icients that are real polynomial in~$\Upsilon$. It follows that every
element of ${\mathcal{R}}$ can be written uniquely in the form
\[ \sum_{p=0}^{k}A_p(\Upsilon)\,\nabla^p\]
for suitable real polynomials $A_p(\Upsilon)$. In our claimed realisation,
such an element is represented by the dif\/ferential operator
\[ \sum_{p=0}^{k}A_p(\tan\theta)\,d^p/d\theta^p\]
and now it suf\/f\/ices to observe (by acting on
$1,\theta,\theta^2,\ldots,\theta^{k}$ near~$0$) that such a dif\/ferential
operator vanishes if and only if all the polynomials $A_p(\Upsilon)$ are zero.
(More precisely, we should restrict the action of such operators to smooth
functions on $(-\pi/2,\pi/2)$ or some other suitable function space.)

Having realised ${\mathcal{R}}$ by dif\/ferential operators, we are reduced to
proving identities amongst these operators. This is accomplished in the
following lemma.\end{proof}

\begin{lemma}\label{tricky} Let $D$ denote the differential operator
$f(\theta)\mapsto (\cos^2\theta)(df(\theta)/d\theta)$ on the circle. Then the
following identities hold. If $k$ is even, then
\[\frac{1}{\cos^k\theta}\frac{d}{d\theta}
\left(D^k\left(\frac{f(\theta)}{\cos^k\theta}\right)\right)
=\left(\frac{d^2}{d\theta^2}+k^2\right)\cdots
\left(\frac{d^2}{d\theta^2}+16\right)
\left(\frac{d^2}{d\theta^2}+4\right)\frac{d}{d\theta}f(\theta)\] and, if $k$ is
odd, then \[\frac{1}{\cos^k\theta}\frac{d}{d\theta}
\left(D^k\left(\frac{f(\theta)}{\cos^k\theta}\right)\right)
=\left(\frac{d^2}{d\theta^2}+k^2\right)\cdots
\left(\frac{d^2}{d\theta^2}+9\right)
\left(\frac{d^2}{d\theta^2}+1\right)f(\theta).\]
\end{lemma}
\begin{proof} Writing $D_{k+1}$ for the operators on the left-hand-sides of the
displays in this lemma,  the following identity is easily verif\/ied
\[
D_{k+3} = \left( \frac{d} { d\theta}
-(k + 2) \tan \theta\right)D_{k+1}\left( \frac{d}{ d \theta} + (k + 2) \tan \theta\right).\]
It follows that, in our realisation of the algebra ${\mathcal{R}}$,
we obtain the expressions on the left-hand-sides of the displays in
the Lemma~\ref{pretricky}.

On the other hand, the operators on the right-hand-sides of the
claimed identities, in the current lemma, are characterised up to scale
as annihilating the functions
\[\cos(k\theta),\quad \sin(k\theta),\quad \cos((k-2)\theta),\quad
\sin((k-2)\theta),\quad \cos((k-4)\theta),\quad \dots.\]
Since all operators have $d^{k+1}/d\theta^{k+1}$ as leading term, it is
therefore suf\/f\/icient to show that the left hand sides of these purported
identities have the same property. Notice that there is an invertible
relationship
\[\cos(m\theta)=2^{m-1}\cos^m\theta+\cdots=T_m(\cos\theta),\]
where $T_m$ is the $m^{\mathrm{th}}$ Chebyshev polynomial of the f\/irst kind
and a similar invertible relationship
\[\sin(m\theta)=(\sin\theta)\big(2^{m-2}\cos^{m-1}\theta+\cdots\big)
=(\sin\theta)U_{m-1}(\cos\theta)\]
where $U_{m-1}$ is the $(m-1)^{\mathrm{st}}$ Chebyshev polynomial of the
second kind. Only the degree of these Chebyshev polynomials concerns us and
it now suf\/f\/ices to show that $D^{k+1}$ annihilates the f\/irst $k+1$ of
\[1,\quad\frac{\sin\theta}{\cos\theta},\quad
\frac{1}{\cos^2\theta},\quad\frac{\sin\theta}{\cos^3\theta},\quad
\frac{1}{\cos^4\theta},\quad\frac{\sin\theta}{\cos^5\theta},\quad
\frac{1}{\cos^6\theta},\quad \frac{\sin\theta}{\cos^7\theta},\quad\dots.\]
Since
\[D\frac{1}{\cos^{2\ell}\theta}=2\ell\frac{\sin\theta}{\cos^{2\ell-1}\theta}
\qquad\mbox{and}\qquad
D\frac{\sin\theta}{\cos^{2\ell-1}\theta}
=(2\ell-1)\frac{1}{\cos^{2(\ell-1)}\theta}
-2(\ell-1)\frac{1}{\cos^{2(\ell-2)}\theta},\]
this follows easily by induction.
\end{proof}
